
%\chapter{Elementos sobre a Economia Política da energia}
%\section{\textit{Machinery Question} e a questão da força-motriz}
%
%\section{Automação \textit{versus} automatização do trabalho}

\setlength{\epigraphwidth}{0.9\textwidth}

\chapter{A literatura sobre o efeito rebound}
\epigraph{Com combustível e fogo, então, quase tudo é fácil. Com a ajuda do fogo, seja no forno de fundição ou na máquina, realizamos, ao longo do último século, sucessivas substituições de um pior por um melhor, de um mais caro por um mais barato, de um velho processo por um novo, os quais impulsionam nossa civilização material. Mas quando esse combustível, nossa energia material, nos faltar, de onde virá a força para realizar feitos iguais ou maiores no futuro? Ninguém pode esperar que, pelo fato de ter realizado muito quando gozava de saúde e vigor físico, fará ainda mais quando suas forças se esvaírem. Contudo, essa é a situação de nossa nação, a menos que a fonte de nossa força seja cuidadosamente preservada ou que se encontre algo melhor que o carvão para substituí-lo e continuar a substituição do pior pelo melhor. Resta saber se o consumo de carvão pode ser reduzido em nosso sistema de livre mercado industrial, ou se, no processo de descoberta, podemos esperar encontrar algum substituto para o carvão. A conclusão imparcial estará longe de ser satisfatória.}{\parencite[p. 72]{Jevons1866}}

\section{Jevons e ``a questão do carvão''}
Ao final do século XIX, em meio à crescente atividade econômica impulsionada pela Revolução Industrial, William Stanley Jevons publicou o livro \textit{The Coal Question} \citeyear{Jevons1866}, em que aborda a crescente preocupação da época sobre os perigos de exaustão das jazidas de carvão da Inglaterra, em meio ao acelerado desenvolvimento industrial ao longo do século. Uma importante conclusão a que Jevons chega, e que posteriormente consolidou-se sob o nome ``paradoxo de Jevons'', é que uma maior eficiência do emprego de carvão induz um \textit{aumento} de seu consumo, oposto à intuição cotidiana que ela induziria uma \textit{diminuição} de seu consumo --- pois, intuitivamente menos energia seria necessária para alcançar um mesmo nível de produto. 
%Em uma famosa passagem, Jevons afirma que ``É uma confusão de ideias supor que o uso econômico de combustível [carvão] seja equivalente à diminuição do [seu] consumo. É o contrário que é verdade.'' \cite[p. 75]{Jevons1866}.

Jevons trata, portanto, do aumento \textit{absoluto} do consumo de carvão enquanto fonte energética, não obstante a redução \textit{relativa} da necessidade de seu consumo para um mesmo nível prévio de produção. Isso pode ser constatado, por exemplo, pelo emprego paulatino de motores com maior eficiência termodinâmica, os quais foram empregados em escalas crescentes de produção (e, portanto, de consumo de matéria-prima e de energia). Ou seja, embora tais ganhos de eficiência \textit{permitam} um consumo menor de carvão para um mesmo nível de produção prévio, o que tende a ocorrer é um aumento do nível de produção \textit{efetivo}. De certa forma, a possibilidade de poupar recursos dá espaço para seu consumo alhures, conduzindo a uma maior atividade econômica e, portanto, maior consumo efetivo de carvão do que previamente demandado.

Jevons faz uma interessante analogia com a ``área de finanças'', em que uma taxa de lucro menor, porém com maior volume de negócios, pode trazer maior receita (``\textit{gross and sometimes even nett revenue}'') do que uma taxa maior de lucros, porém com volume menor de negócios. Da mesma forma, conclui ele, uma maior eficiência do consumo de energia permite com que um consumo \textit{maior} de energia traga maior receita. Ou seja, como conclui, ``É justamente a economia de seu uso que leva ao seu consumo extensivo.'' \cite[p. 76]{Jevons1866}. 

Um ponto importante de se destacar, e muito frequentemente olvidado, é que Jevons faz uma separação explícita entre o uso \textit{doméstico} do carvão e o seu uso \textit{industrial/manufatureiro}. Quanto ao primeiro, afirma:
\begin{quote}
    Este [consumo doméstico do carvão], sem dúvida, pode ser reduzido sem maiores prejuízos além da diminuição do conforto em nossos lares e de uma leve alteração em nossos hábitos nacionais consolidados. O carvão assim economizado seria, em sua maior parte, armazenado para uso das gerações futuras. Mas mesmo que nossa população pudesse ser induzida a se abster do prazer de uma boa lareira, a economia [\textit{saving}] obtida não ultrapassaria cerca de um terço do consumo total de carvão; o consumo doméstico é, em média, de cerca de uma tonelada por ano, por habitante. \textit{Dos outros dois terços, quase um terço é utilizado em nossas siderúrgicas; e o restante em nossas fábricas, fornos e oficinas mecânicas em geral.} \cite[p. 75, grifo meu]{Jevons1866}
\end{quote}

Por outro lado, o trecho mais comumente citado de Jevons junta dois enunciados em seções diferentes de \textit{The Coal Question}: juntam um enunciado no capítulo I, em que ele diz que
\begin{quote}
    Fica demonstrado [`nos capítulos vi., vii. e viii.'] que a tendência constante das descobertas [quanto a usos do carvão] é tornar o carvão um agente cada vez mais eficiente, enquanto não haja probabilidade de que, quando nosso carvão se esgotar, surja um substituto mais potente. \textit{Tampouco o uso econômico do carvão reduzirá seu consumo. Pelo contrário, a economia torna o emprego do carvão mais lucrativo e, assim, a demanda atual por carvão aumenta}, e a vantagem se inclina mais fortemente para o lado daqueles que, no futuro, terão os suprimentos mais baratos. \parencite[p. 17, grifo meu]{Jevons1866}
\end{quote}
e um enunciado no capítulo VII, onde ele resume explicitamente seus resultados elaborados nas seções anteriores:
\begin{quote}
    Mas a economia do carvão nas manufaturas é uma questão diferente. \textit{É uma completa confusão de ideias supor que o uso econômico de combustível seja equivalente a um consumo reduzido. O contrário é que é verdade.}'' \parencite[p. 75, grifo meu]{Jevons1866}
\end{quote}

São as frases destacadas em itálico nestas duas citações que são explicitamente referenciadas juntas, em particular, por Brookes (\citeyear*[p. 319]{Brookes1990}; \citeyear*[p. 357, n. 8]{Brookes2000})\footnote{Em ambos, Brookes cita-o da mesma forma: ``\textit{Nor will the economical use of coal reduce its consumption. On the contrary, economy renders the employment of coal more profitable and thus the present demand for coal is increased...} [perceba a `inócua' elipse]. \textit{It is wholly a confusion of ideas to suppose that the economical use of fuel is equivalent to diminished consumption. The very contrary is the truth}.''.}, ofuscando não só que há 7 capítulos entre ambas, como também que as últimas duas frases tratam-se do emprego \textit{manufatureiro} do carvão! Tal conflação do consumo produtivo e do improdutivo é um importante ponto de atenção na literatura neoclássica deste fenômeno, como será melhor visto adiante.

% Ao final do século XIX, William Stanley Jevons escreve ``\textit{The Coal Question}'' \citeyear{Jevons1866}, em que discute sobre os perigos de exaustão das jazidas de carvão da Inglaterra, em meio ao acelerado desenvolvimento industrial ao longo do século. No que tange ao principal consumo de carvão --- a fonte de energia prevalente à época ---, Jevons destaca que o consumo doméstico \textit{não} desempenha peso significativo:

% Em vez disso, dá destaque sobre o caráter do ``consumo das manufaturas'' --- i.e. da indústria, consumo produtivo ---, sendo esta a sua famosa citação à qual autores usualmente aludem:
% \begin{quote}
%     ``But the economy of coal \textit{in manufactures} is a different matter. It is wholly a confusion of ideas to suppose that the economical use of fuel is equivalent to a diminished consumption. The very contrary is the truth.'' (Ibid.; grifo meu)
% \end{quote}

% Este é o cerne do que se convencionou chamar de ``paradoxo de Jevons'': uma maior eficiência do emprego da energia induz um \textit{aumento} do consumo (a um nível da economia como um todo), oposto à intuição cotidiana que ela induziria uma \textit{diminuição} de seu consumo (pois menos energia seria necessária para alcançar um mesmo nível de produto).\footnote{“\textit{And if such is not always the result within a single branch, it must be remembered that the progress of any branch of manufacture excites a new activity in most other branches, and leads indirectly, if not directly, to increased in roads upon our seams of coal.}” \cite[p. 76]{Jevons2018}.}

% % Jevons faz uma analogia com a ``área de finanças'', em que uma taxa de lucro menor, porém com maior volume de negócios, pode trazer maior receita (``\textit{gross and sometimes even nett revenue}'') do que uma taxa maior de lucros, mas com volume menor de negócios; da mesma forma, uma maior eficiência do consumo (produtivo) de energia permite com que um consumo (produtivo) \textit{maior} de energia traga maior receita. Ou seja, como conclui, "\textit{It is the very economy of its use which leads to its extensive consumption.}" \cite[p. 76]{Jevons2018}. 

% Jevons trata, portanto, do aumento \textit{absoluto} do consumo de carvão enquanto fonte energética (produtiva ou improdutivamente), não obstante sua redução \textit{relativa} de consumo (o qual torna-se, de fato, ``\textit{mais eficiente}''). Isso pode ser visto, por exemplo, pelo emprego paulatino de motores com maior eficiência (termodinâmica), suplementação da força-motriz do carvão com a força-motriz a vapor (o foco da "segunda Revolução Industrial"), etc. 


\section{Os debates neoclássicos e o postulado de Khazzoom-Brookes}
A discussão sobre o ``paradoxo'' de Jevons obscureceu-se ante o advento da exploração do petróleo na virada para o século XX e, eventualmente, da energia nuclear, através dos quais os limites energéticos --- que outrora estavam à espreita das preocupações dos cientistas e dos capitalistas --- estavam agora como que superados. Gás, petróleo e eletricidade, que eram empregados como fontes de energia em iluminação urbana no final do século XIX e início do século XX, tinham uma queda tendencial em seus preços, o que instigou seu emprego como fontes energéticas em outras indústrias, das quais destaca-se a indústria automotiva de motores de combustão interna, movidos a gasolina \parencite[p. 8-9]{Fouquet2009}. 


Tais discussões somente voltam a assumir uma relevância maior quando dos choques do petróleo na década de 1970, em que os aumentos extremos dos preços dos barris de petróleo, levados a cabo pelos países da OPEP, induziram uma recessão mundial e requereram medidas estatais para estabilização das economias em âmbito nacional. É neste contexto que \textcite{Khazzoom1980} escreve seu texto seminal sobre os padrões de eficiência que foram estatalmente estabelecidos, nos Estados Unidos, em particular no que tange ao uso energético de carros e eletrodomésticos. Não é difícil lembrar das famosas fotos de filas de carros em frente a postos de gasolina, de pessoas tentando garantir o mínimo de combustível que conseguissem --- eram tais preocupações que estavam candentes no imaginário público ao final da década de 1970.

Neste artigo, Khazzoom elabora sobre uma noção central na literatura sobre o efeito \textit{rebound} que o sucedeu: ganhos de eficiência se refletem em mudanças de ``preço implícito''. Para ilustrá-lo, ele utiliza-se do imaginário cotidiano: um indivíduo que compre um forno com o dobro de eficiência --- com relação a um padrão pré-ganho de eficiência --- fará com que ``o `preço efetivo' do combustível [empregado] caia pela metade: ele custará agora metade do que custava para manter a fornalha operando a cada hora'', da mesma forma como um carro que faça, por exemplo, 30 quilômetros por litro de gasolina ``dá na mesma'' (\textit{ceteris paribus}) que um carro que faça 10 quilômetros por litro, porém com preços de gasolina $1/3$ menores \parencite[p. 22]{Khazzoom1980}. %(Nestes exemplos, há dois pressupostos importantes: que o preço do combustível permanece o mesmo mediante tais ganhos de eficiência, e que o usufruto destes bens também permanece igual.)

Como forma de desnaturalizar tais intuições, Khazzoom menciona um exemplo típico de Microeconomia: intuitivamente pensa-se que, caso 100 indivíduos fossem empregados em algum processo produtivo, então tal quantidade ``cairia pela metade'' mediante a duplicação da produtividade do trabalho. Porém, isso somente é o caso caso a demanda por tais produtos \textit{permaneça a mesma}, o que certamente não é o que acontece na prática: ``a demanda não permanece em seu nível prévio, mas tende a aumentar'' \parencite[p. 23]{Khazzoom1980}, em cujo caso empregar-se-á não todos os mesmos 100 indivíduos, é claro, mas certamente \textit{mais do que a metade}.
%\footnote{Note-se que, aqui, Khazzoom emprega um exemplo \textit{do setor produtivo} para desmistificar a aparência de fenômenos \textit{domésticos/improdutivos}. Como será destacado posteriormente, isso é interessante uma inversão dialética: embora a produção (ou ``consumo produtivo'') possua prioridade ontológica com relação ao consumo (``improdutivo'') --- porquanto não se pode consumir aquilo que não foi produzido previamente \parencite{Marx2011} ---, é justamente a esfera do consumo que nos aparece mais imediatamente no cotidiano, dando a impressão que fenômenos econômicos na esfera da produção são análogos àqueles da esfera da circulação e do consumo.}

A novidade neste artigo de Khazzoom é de formalizar, dentro do próprio panorama neoclássico, o argumento de que um aumento de eficiência energética induz um aumento do consumo. Sua análise é feita com argumentos microeconômicos, concluindo com o enunciado de que 
\begin{quote}
    A elasticidade da demanda de energia para um uso final em relação à eficiência do seu próprio aparelho é igual ao negativo da soma de 1 mais a elasticidade da demanda de energia para esse uso final em relação ao preço da energia. \parencite[p. 31]{Khazzoom1980}
\end{quote}

Ou, formalmente, chamando a elasticidade-eficiência da demanda energética de $\epsilon_{Ef}$ e a elasticidade-preço dessa demanda de $\epsilon_P$, o resultado de Khazzoom é da forma
$$
\epsilon_{Ef} = -(1+\epsilon_P)
$$

Ou seja, dizer que um aumento na eficiência energética induz uma queda aproximadamente equivalente na demanda é o mesmo que dizer que $\epsilon_{Ef} \approx -1$, o que é o mesmo (pela derivação de Khazzoom) que dizer que $\epsilon_P \approx 0$, i.e. que a demanda energética --- ou seja, tanto o estoque de aparelhos que empregam energia quanto a taxa de utilização deles --- \textit{não} seja afetada pelo preço da energia, o que contrasta flagrantemente com o que é verificado empiricamente \parencite[p. 39]{Khazzoom1980}.

% \begin{itemize}
%     \item Mudanças de eficiência refletem-se em mudanças de ``\textit{preço implícito}''
%     \item Por exemplo, um forno com eficiência duas vezes melhor fará com que ``o `preço efetivo' do combustível caia pela metade: ele custará agora metade do que custava para manter a fornalha operando a cada hora'' \parencite[p. 22]{Khazzoom1980}
%     \item Ou seja, sua demanda por combustível com um forno mais eficiente --- em, digamos, uma hora --- é metade do que era com o forno menos eficiente, \textit{mantido o mesmo uso que antes}. Ou seja, de certa forma, o preço implícito do combustível --- no que tange a este forno --- diminuiu pela metade
%     \item \textit{Mutatis mutandis} para um carro mais eficiente: se ele faz o dobro da quilometragem do que um carro menos eficiente, então o seu ``custo por quilômetro'' cairá pela metade, pois cada quilômetro demandará a metade do combustível de antes
%     \item Nisso está pressuposto duas coisas: que o preço do combustível permanece o mesmo, e que o seu uso permanecerá no mesmo nível que antes
%     \item Khazzoom faz um paralelo com o aumento da produtividade do trabalho: caso a produtividade dobre, Khazzoom diz que, de acordo com ``a literatura no começo e no meio do século XIX'', ``metade das pessoas que trabalhavam na produção até ali tornar-se-iam redundantes, e ficariam desempregadas'' \parencite[p. 23]{Khazzoom1980}
%     \item Porém, este argumento pressupõe que a demanda \textit{permaneceria a mesma}, o que não é o caso: ``a demanda não permanece em seu nível prévio, mas tende a aumentar'' (ibid.).
%     \item Ou seja, ``no que diz respeito ao produtor'', a situação é análoga à do consumidor --- em verdade, o que acontece é o contrário: é o caso do produtor que é assim, e estamos extrapolando a analogia para o consumidor! 

Entrementes, ao longo da década de 1980, outro tema estava ganhando tração dentro da área da Economia da Energia, em particular no que tangia a políticas estatais/públicas: ``quais medidas seriam mais eficientes quanto à mitigação do aquecimento global\footnote{Ou também aludido como ``resfriamento global'' (\textit{global cooling}) em países do Norte Global.}?''. Uma visão à época, representada e.g. por \textcite{Keepin1988,Keepin1989}, argumentava a favor de medidas estatais que buscassem \textit{aumentar a eficiência energética nacional} como ``[a] resposta mais rápida, menos dispendiosa e, sobretudo, mais eficaz ao aquecimento global provocado pela queima de combustíveis fósseis'' \parencite[p. 554]{Keepin1988}.\footnote{É justo destacar que o trabalho de \textcite{Keepin1988} buscou uma análise comparativa entre duas estratégias contra o aquecimento global: a transição de combustíveis fósseis (altos emissores de GEEs) para energia nuclear (``... a geração de eletricidade nuclear não acarreta emissões diretas de $\texttt{CO}_2$ ou de quaisquer outros gases de efeito estufa'', p. 540) \textit{versus} o aumento da eficiência do uso energético, em cujo caso concluem que esta última é bem mais custo-benefício e realisticamente implementável. Mesmo assim, não se pode negar que o propósito do artigo é inegavelmente normativo: ele não só argumenta que estas políticas são as mais eficientes, mas também de que é cabível que elas \textit{sejam levadas a cabo}.} Ou seja, a questão de ganhos de eficiência energética assumia agora não só o papel de mitigar a dependência econômica de combustíveis importados, como também o papel de política doméstica de mitigação do efeito estufa.

Levando adiante o argumento de Khazzoom, \textcite{Brookes1990,Brookes1990a}\footnote{Brookes publicou \textit{The greenhouse effect: the fallacies in the energy efficiency solution} em março de 1990 \parencite{Brookes1990a} e \textit{Energy Efficiency and The Greenhouse Effect} em dezembro de 1990 \parencite{Brookes1990}. Como este último não só engloba os argumentos do primeiro, como elabora-os de forma mais detida, optei por empregá-lo como fio condutor de sua discussão à época.} argumenta contra as propostas destes economistas --- os quais passaram a ser alcunhados de ``conservacionistas'' (de energia). Brookes elabora a mesma ideia que Khazzoom, partindo da conclusão de Jevons supracitada. O cerne de seu argumento é que ``um recurso encontrar-se em um mundo de uso mais eficiente é o mesmo que ele desfrutar de uma redução em seu preço implícito, com tudo o que isso implica para a demanda''; isto é, um mesmo nível de usufruto de algum recurso que se torna mais eficiente é equivalente a que o preço (nominal) deste mesmo usufruto tivesse caído \parencite[p. 319]{Brookes1990}. 
% \begin{quote}
%     um recurso encontrar-se em um mundo de uso mais eficiente é o mesmo que ele desfrutar de uma redução em seu preço implícito, com tudo o que isso implica para a demanda. Isso significa que, se consumidores (e produtores) descobrirem que obtêm mais benefícios com um determinado nível de gasto em um recurso (energia, neste caso), [então] isso terá o mesmo efeito em suas decisões como se o preço tivesse caído: o recurso em questão torna-se um item relativamente mais atraente em comparação com outros em sua lista de compras. \parencite[p. 319]{Brookes1990} 
% \end{quote}

É por isso que Brookes critica que tais medidas de eficiência energética ``não são mais eficientes como soluções ao efeito estufa do que como soluções a problemas de oferta energética [\textit{energy supply problems}]'' \parencite[p. 320]{Brookes1990}. Ou seja, propõe que tais medidas sejam vistas ``pelo que elas são'': como ``auxílios ao crescimento econômico'' (ibid.), e que, \textit{neste âmbito}, elas podem, aí sim, ser úteis no que tange ao aquecimento global, embora indiretamente: ao permitirem um maior nível de produto nacional, fornece-se os meios pelos quais futuras gerações poderão ``enfrentar os problemas que elas realmente poderão enfrentar --- os quais podem ser muito diferentes daqueles que nós, com nossa visão essencialmente míope do futuro, podemos projetar erroneamente sobre eles'' \parencite[p. 319]{Brookes1990}. Brookes assevera que é preciso considerar ``muito criticamente'' o apoio a tais campanhas nacionais de aprimoramento de eficiência energética, vistas por governos (p. ex. do Reino Unido, no contexto de Brookes) como fatores importantes para ``reduzir a dependência em combustível importado e para lidar com a ameaça do aquecimento global'', pois --- e, quanto a isso, Brookes não mede palavras --- ``se estiverem errados, recursos econômicos valiosos poderão ser sacrificados a falsos deuses'' \parencite[p. 321]{Brookes1990}. 

À guisa de crítica, Brookes dá exemplos de como os principais índices de eficiência energética, empregados pelos ``conservacionistas'', são enviesados \parencite[p. 321]{Brookes1990}. O coeficiente de energia (\textit{energy coefficient}) é definido como ``a razão entre a taxas de crescimento da energia e a taxa de crescimento econômico'', ou seja,
$$
C_E \equiv \frac{\dot{E}}{\dot{Y}} = \frac{\frac{dE}{dt}}{E} \cdot \frac{Y}{\frac{dY}{dt}}
$$
e a razão energética (\textit{energy ratio}), definida como ``o consumo de energia por unidade de produto econômico'', ou seja,
$$
R_E \equiv \frac{E}{Y}
$$ 

Pode-se demonstrar que\footnote{Fazendo a derivada de $R_E$: $\frac{d}{dt}\left(E/Y\right) = \frac{dE}{dt}\frac{1}{Y} - \frac{E}{Y} \left(\frac{dY}{dt}\frac{1}{Y}\right)$. Dividindo por $R_E$ para obter a taxa de crescimento $\dot{R_E}$: $\frac{1}{E/Y} \frac{d}{dt}\left(E/Y\right) = \dot{\left(E/Y\right)} \equiv \dot{R_E} = \frac{dE}{dt}\frac{1}{E} - \frac{dY}{dt}\frac{1}{Y} = \dot{E} - \dot{Y}$. Dividindo por $\dot{Y}$ e rearranjando: $\dot{E}/\dot{Y} \equiv C_E = \left(1 + \frac{\dot{R_E}}{\dot{Y}}\right)$.}
$$
C_E = \left(1 + \frac{\dot{R_E}}{\dot{Y}}\right)
$$
em cujo caso é possível que ocorra que o coeficiente $C_E$ esteja em queda --- i.e. \textit{``aponte aumentos energéticos''} --- mesmo em casos em que não haja ganhos de eficiência \textit{de facto}. Brookes dá um exemplo numérico empírico \parencite[p. 321-2]{Brookes1990}: ele aponta uma queda secular no consumo energético por unidade de produto de $1\%$ por ano --- i.e. $\dot{R_E} = -1\%$ ---, em cujo caso basta que $\dot{Y} \leq 1\%\, a.a.$ para que $C_E \geq 0$. Ou seja, este coeficiente é enviesado no que tange à medição de ganhos de eficiência em uma economia, porquanto pode dar ``falsos positivos'' em épocas de recessão e/ou de choques em \textit{supply chains}, como foi o caso em 1973 e 1979. 

Quanto ao segundo termo --- a razão energética $R_E$ ---, Brookes aponta para resultados empíricos de Schurr, nos EUA, que apontam que ``a substituição de capital e trabalho por energia leva ao crescimento da produtividade multifatorial'' \parencite[p. 325]{Brookes1990}. Ou seja, 
\begin{quote}
    Se a produtividade multifatorial melhorar (talvez devido a práticas de trabalho mais eficientes ou à substituição de instalações menos produtivas por instalações mais produtivas), cria-se a impressão de melhoria da produtividade energética, mesmo que não tenha havido nenhuma mudança nas práticas de consumo de energia. \parencite[p. 326]{Brookes1990}
\end{quote}
em cujo caso o denominador de $E/Y$ pode aumentar sem que o consumo (\textit{absoluto}) de $E$ mude. Brookes conclui, portanto, que ambos estes termos são insuficientes para mensurar a questão de ganhos de eficiência energética, sendo particularmente perniciosos nas discussões de políticas públicas sobre mitigação do efeito estufa.

Em suma, Brookes remete à noção mais fundamental da Economia (Ortodoxa): que o homem econômico (``\textit{economic man}'') 
\begin{quote}
    busca alocar todos os recursos para maximizar sua satisfação econômica (otimizar seus retornos como trabalhador e poupador), sujeito a todas as restrições que o cercam. Qualquer melhoria na produtividade energética que advenha de suas ações seria \textit{incidental e subordinada ao seu objetivo principal}. \parencite[p. 328, grifo meu]{Brookes1990}
\end{quote}
Nesse ínterim, o foco compenetrado nos ganhos de eficiência --- e ainda mais na eficiência \textit{energética} --- teria mais chances de ser sub-ótimo do que o contrário --- e, mesmo que o seja, por que não buscar justamente aquilo que se deveria buscar, i.e. ``maximizar o bem-estar sujeito a todas as restrições circundantes'' (ibid.)?
%\footnote{\textcite{Grubb1990} argumenta que, ``[s]e a eficiência energética é o \textit{objetivo}, invés de os \textit{meios}, então a força motriz [do aumento da demanda energética] é completamente diferente --- assim como as implicações macroeconômicas'' (p. 783). }


% Brookes:
% \begin{itemize}
% %    \item ``Isso significa que, se consumidores (e produtores) descobrirem que obtêm mais benefícios com um determinado nível de gasto em um recurso (energia, neste caso), [então] isso terá o mesmo efeito em suas decisões como se o preço tivesse caído: o recurso em questão torna-se um item relativamente mais atraente em comparação com outros em sua lista de compras.'' \parencite[p. 319]{Brookes1990} [Note-se que aqui está mais explícito: este efeito afeta consumidores \textit{e também produtores}!]  
% %    \item Brookes critica que tais medidas de eficiência energética --- embora sejam defendidas pelos ``conservacionistas'' como, por exemplo, ``[a] resposta mais rápida, menos dispendiosa e, sobretudo, mais eficaz ao aquecimento global provocado pela queima de combustíveis fósseis é reduzir a emissão de CO2'' \parencite[p. 554]{Keepin1988} ---, ``não são mais eficientes como soluções ao efeito estufa do que como soluções a problemas de oferta energética [\textit{energy supply problems}]'' \parencite[p. 320]{Brookes1990}
% %    \item Brookes argumenta, neste ínterim, que tais medidas sejam vistas ``pelo que são'': como ``auxílios ao crescimento econômico'' (ibid.), e que, \textit{neste âmbito}, elas podem, aí sim, ser úteis no que tange ao aquecimento global, embora indiretamente: ao permitirem um maior nível de produto nacional, fornece-se os meios pelos quais futuras gerações poderão ``enfrentar os problemas que elas realmente poderão enfrentar --- os quais podem ser muito diferentes daqueles que nós, com nossa visão essencialmente míope do futuro, podemos projetar erroneamente sobre eles'' \parencite[p. 319]{Brookes1990}. Brookes assevera que é preciso considerar ``muito criticamente'' o apoio a tais campanhas nacionais de aprimoramento de eficiência energética, vistas por governos (do Reino Unido, à época de Brookes) como fatores importantes para ``reduzir a dependência em combustível importado e para lidar com a ameaça do aquecimento global'', pois --- e, quanto a isso, Brookes não mede palavras --- ``se estiverem errados, recursos econômicos valiosos poderão ser sacrificados a falsos deuses'' \parencite[p. 321]{Brookes1990}.
%     \item Falácias apontadas por Brookes: 
%     \item falácia da composição: ``acreditar que as reduções no consumo de energia no nível \textit{microeconômico} de aplicações individuais [\textit{individual applications}] levarão automaticamente a reduções proporcionais no nível \textit{macroeconômico}.'' \parencite[p. 321, grifo meu]{Brookes2000}
%     \item Falácia de ``massa de atividades dependentes de energia'' (``\textit{lump-of-energy-dependent-activity fallacy}''): aludindo às falácias de ``massas de trabalho/capital'' (\textit{lump of labour/capital fallacies}), aponta quanto a pressupor-se que a quantidade de atividades dependentes de energia permanecerão intactas mediante tais alterações de preço implícito da energia e consequentes reconfigurações de fatores de produção 
%     \item O que ignoram é que 1) tais quedas de preço implícito aumentam a demanda tanto quanto uma queda de preço nominal, e que 2) as economias de energia advindas de tais ganhos de eficiência ``precisam achar alguma destinação [\textit{outlet}]. Nas sociedades industriais modernas, é provável que esta destinação esteja nos bens e serviços que requerem energia em sua produção, se não também em seu uso.'' \parencite[p. 321]{Brookes2000}
% \end{itemize}
